\documentclass[12pt]{article}
\usepackage[utf8]{inputenc}
\usepackage[T1]{fontenc}
\usepackage{hyperref}
\title{Suau Language Resources for K-12 Curriculum Development}
\author{Compiled by Claude Sonnet 5.5 (claude-sonnet-5-5) for Andr\'as Kornai}
\date{}
\begin{document}
\maketitle

\section{Introduction}
Suau (ISO 639-3: swp; Glottocode suau1242) is an Oceanic (Austronesian) language of Milne Bay Province, Papua New Guinea. Sources give about 6,800 to 7,800 speakers. Dialects named in one source include Bona Bona, Bonarua, Dahuni, Lausaha (Daui, Fife Bay), Leileiafa and Sinaki (Gadaisu). Joshua Project notes that younger people may shift towards Tok Pisin or English.

\section{Primary Educational Resources}
The searches found no school primers or textbooks. The new dictionary below is the nearest thing to a teaching resource.

\section{Linguistic and Reference Materials}
\begin{itemize}
\item News report on the launch of the Suau dictionary \emph{Ta alina Suau} (``Let's Speak Suau''), by Chris Abel with villagers, after about 34 years of work (EMTV): \url{https://emtv.com.pg/new-dictionary-for-suau-people}
\item Verbix page on Suau, listing dialects and a speaker count: \url{https://docs.verbix.com/Languages/Suau}
\item Wikipedia entry for Suau: \url{https://wikipedia.com/wiki/ISO_639:swp}
\end{itemize}

\section{Religious and Community Resources}
\begin{itemize}
\item Suau Bible text \emph{Ribaharihariuna} (New Testament and parts of the Old Testament, Bible Society of Papua New Guinea; the search result says it is under a CC BY-NC-ND 4.0 licence): \url{https://ebible.org/pdf/swp/swp_all.pdf}
\item Joshua Project profile, Suau-Daui (Fife Bay): \url{https://joshuaproject.net/index.php/people_groups/15097}
\item Joshua Project profile, Suau (Sinaki): \url{https://joshuaproject.net/index.php/people_groups/14189}
\end{itemize}

\section{Digital and Audio Resources}
The searches found no audio, video or software resources for Suau. The Bible is available as a PDF.

\section{Curriculum Development Context}
Suau has a published dictionary and a scripture translation, and so an established written form. The searches found no evidence of school use. A curriculum could draw on the dictionary and the Bible text, subject to the dictionary's availability and permission from its authors and community.
\end{document}